\documentclass[11pt]{report}

\usepackage[margin=1in]{geometry}
\usepackage{booktabs}
\usepackage{graphicx}
\usepackage{amsmath}
\usepackage{siunitx}
\sisetup{group-separator={,}, detect-weight=true}
\usepackage{hyperref}
\hypersetup{colorlinks=true, linkcolor=black, urlcolor=blue}

\title{Spara: Partition-Aware Distributed GNN Training\\[4pt]
       \large Dissertation Defense --- Experimental Results}
\author{Sandeep Polisetty}
\date{\today}

\begin{document}
\maketitle
\tableofcontents

\chapter{Experimental Evaluation}

Spara (\texttt{gsplit}) is an efficient distributed GNN training system built on top
of the Deep Graph Library (DGL). It introduces \emph{partition-aware sampling} and
\emph{feature loading} to reduce cross-GPU communication during multi-GPU GNN
training: each GPU owns a METIS partition (weighted by sampling frequency), and
sampling, feature caching, and the model forward/backward exchange only boundary data
across GPUs via NCCL. We are additionally exploring \emph{hybrid parallelism}: for
certain hidden sizes and depths the split communication cost becomes too large, and a
hybrid of data- and split-parallelism should do better. To choose the best strategy
\emph{without} building every variant, we use a cost-model simulator; this chapter
reports the measured system performance (Section~\ref{sec:measured}) and how faithfully
the simulator reproduces it (Section~\ref{sec:simulator}).

\paragraph{Setup.}
All experiments ran on \textbf{RunPod}, on a node with \textbf{2$\times$ NVIDIA
A100-SXM4-80GB} GPUs (NVLink/NVSwitch). Software: \textbf{CUDA 12.4},
\textbf{PyTorch 2.4 (cu124)}, and a custom \textbf{DGL} build with the Spara
C++/CUDA operators. We run the \texttt{split} (Spara) system on OGB
node-classification graphs with \textbf{fanout 15 per layer}, hidden size 256, batch
size 1024, real OGB features, on \textbf{2 GPUs} (\texttt{world\_size}$=2$, 2-partition
map \texttt{\_w2\_ndst\_efreq\_xbal}, UVA mode), sweeping depth over $\{3,4,5\}$ and
models \{GraphSAGE, GAT\}. All times are per-epoch averages.

\section{Measured Spara Performance}
\label{sec:measured}

Table~\ref{tab:measured} reports, per configuration, the time spent in the three
phases: \textbf{sample} (neighbor sampling + the partition-aware neighbor exchange),
\textbf{load} (feature gather to the GPU), and \textbf{train} (forward $+$ backward).
Test accuracy confirms the pipeline trains correctly on real features.

\begin{table}[h]
\centering
\caption{Measured Spara (\texttt{split}) per-epoch breakdown. Fanout is 15 per layer;
hidden size 256; 2$\times$A100-SXM.}
\label{tab:measured}
\begin{tabular}{l l r r r r r r r}
\toprule
Graph & Model & Layers & Fanout & Hidden & Sample (s) & Load (s) & Train (s) & Acc (\%) \\
\midrule
ogbn-arxiv      & SAGE & 3 & 15 & 256 & 0.46 & 0.03 & 0.86 & 56.9 \\
ogbn-arxiv      & SAGE & 4 & 15 & 256 & 0.53 & 0.04 & 1.19 & 54.9 \\
ogbn-arxiv      & SAGE & 5 & 15 & 256 & 0.61 & 0.02 & 1.54 & 56.3 \\
\midrule
ogbn-products   & SAGE & 3 & 15 & 256 & 1.26 & 0.09 & 2.05 & 86.5 \\
ogbn-products   & SAGE & 4 & 15 & 256 & 4.32 & 0.03 & 8.11 & 86.7 \\
ogbn-products   & SAGE & 5 & 15 & 256 & 8.61 & 0.22 & 22.65 & 87.2 \\
\midrule
ogbn-products   & GAT  & 3 & 15 & 256 & 1.30 & 0.09 & 5.80 & 85.8 \\
ogbn-products   & GAT  & 4 & 15 & 256 & 4.16 & 0.01 & 15.66 & 86.5 \\
ogbn-products   & GAT  & 5 & 15 & 256 & 8.76 & 0.24 & 44.85 & 86.3 \\
\midrule
ogbn-papers100M & SAGE & 3 & 15 & 256 & 10.61 & 0.11 & 9.33 & 45.5 \\
ogbn-papers100M & SAGE & 4 & 15 & 256 & 26.84 & 4.42 & 19.49 & 46.1 \\
ogbn-papers100M & SAGE & 5 & 15 & 256 & 48.36 & 31.66 & 47.22 & 45.3 \\
\bottomrule
\end{tabular}
\end{table}

\paragraph{Where the bottleneck is (and isn't).}
Neighbor sampling is \emph{not} the primary target for parallelism: the local sampling
work depends only on topology and fanout, is \textbf{independent of hidden size}, and
across our runs is relatively insensitive to depth. It is also
\emph{strategy-independent}---incurred identically under data, split, or hybrid
parallelism---so changing the parallelism strategy cannot reduce it. (In the
\texttt{split} measurements the ``Sample'' column additionally bundles the
partition-aware neighbor exchange over NCCL, i.e.\ communication, which is why it is
large.) The costs that actually drive the parallelism trade-off are \textbf{feature
loading} and \textbf{cross-partition communication}, which grow with hidden size and
depth---these are what the simulator models and what hybrid parallelism aims to reduce.
On papers100M the load cost jumps at depth~5 as the working set exceeds the on-GPU
feature cache and more features are gathered from host memory over UVA; and GAT's
heavier per-layer compute roughly doubles train time vs.\ SAGE at equal depth.

\section{Simulator Fidelity}
\label{sec:simulator}

To pick a hybrid-parallelism strategy before building it, Spara includes a cost-model
simulator (\texttt{get\_best\_strategy.py} $\rightarrow$
\texttt{simulation/hybrid\_dp\_sim.py}). For a given graph, model, depth and hidden
size it sweeps the \emph{split-point} $k$ (the layer from which split parallelism is
used): $k=0$ is full split parallelism---identical to the measured \texttt{split}
system---and $k=L$ is full data parallelism. For each $k$ it predicts the per-epoch
\textbf{computation} cost (model-layer compute, timed on GPU), \textbf{communication}
cost (cross-partition feature volume $\div$ interconnect bandwidth), and
\textbf{data-loading} cost (frontier feature transfer $\div$ PCIe). The simulator writes
results through the same logging path as the measured runs, so the two are directly
comparable.

To validate the simulator we compare its $k=0$ prediction against the measured numbers
from Section~\ref{sec:measured}: predicted \emph{computation} vs.\ measured
\emph{train}, and predicted \emph{data-loading} vs.\ measured \emph{load}. (Sampling
time is deliberately not modelled---see above.)

\begin{table}[h]
\centering
\caption{Simulator ($k{=}0$, full split) vs.\ measured, per-epoch. \emph{Simulator run
in progress---values to be inserted once \texttt{run\_simulation.sh} completes.}}
\label{tab:sim}
\begin{tabular}{l l r r r r r}
\toprule
 & & \multicolumn{2}{c}{Train / Compute (s)} & \multicolumn{2}{c}{Load (s)} & Comm (s) \\
\cmidrule(lr){3-4}\cmidrule(lr){5-6}
Graph & Layers & Measured & Sim & Measured & Sim & Sim \\
\midrule
ogbn-products & 3 & 2.05  & --- & 0.09 & --- & --- \\
ogbn-products & 4 & 8.11  & --- & 0.03 & --- & --- \\
ogbn-products & 5 & 22.65 & --- & 0.22 & --- & --- \\
\midrule
ogbn-arxiv    & 3 & 0.86  & --- & 0.03 & --- & --- \\
ogbn-arxiv    & 4 & 1.19  & --- & 0.04 & --- & --- \\
ogbn-arxiv    & 5 & 1.54  & --- & 0.02 & --- & --- \\
\bottomrule
\end{tabular}
\end{table}

\paragraph{Calibration notes.}
Two model choices affect how closely the prediction matches hardware and are being
validated: (i) the interconnect model (NVLink vs.\ PCIe bandwidth) must match the
A100-SXM NVLink fabric; (ii) the mapping from sampled hops to convolution layers in the
compute-cost loop. These are being checked against the measured numbers above.

\end{document}
